\section{Robustness Study}
\label{sec:robustness}

The comparison of Section~\ref{sec:sil-comparison} holds for the nominal
plant --- the very model both controllers were designed on. This section
asks how much of that picture survives when the plant is no longer the
design model. The protocol is deliberately strict: \emph{only the plant
changes}. Both controllers keep the gains and trim commands of the nominal
design, and every run repeats the combined mission of
Section~\ref{sec:sil-comparison} ($\VT$ $12\!\to\!15.2$\,m/s at 5\,s,
heading doublet $\pm5^\circ$ at 50/100/150\,s, altitude doublet
$\pm5$\,m at 80/160/240\,s, 300\,s total) --- the same combined
mission, and the same per-axis reference set, used throughout
Section~\ref{sec:sil}. Four scenarios are considered,
spanning plausible-to-severe model uncertainty and an exogenous
disturbance: (A)~a severe model mismatch, replacing the complete
aerodynamic set by an independent identification of the same airframe;
(B)~the longitudinal half of that mismatch only, to isolate the origin of
the failure observed in~(A); (C)~a plausible identification uncertainty,
multiplying all 23 stability and control derivatives by 1.10; and (D)~wind
and gusts applied through the native disturbance inputs
of~\eqref{eq:inputvec}. Table~\ref{tab:robustness} summarizes the tracking
metrics; the degradation columns normalize each scenario by the same
controller's nominal performance, so that the two laws are compared
against themselves rather than against each other.

\begin{table*}[!t]
\caption{Robustness campaign: RMS tracking errors ($t>12$\,s) and
degradation relative to the nominal scenario (RMS scenario / RMS nominal).
Gains and trim always from the nominal design.}
\label{tab:robustness}
\centering
\begin{tabular}{llcccccc}
\toprule
 & & \multicolumn{3}{c}{RMS error} & \multicolumn{3}{c}{Degradation ($\times$ nominal)} \\
Scenario & Ctrl. & $\VT$ [m/s] & $\psi$ [$^\circ$] & $h$ [m] & $\VT$ & $\psi$ & $h$ \\
\midrule
Nominal & PID & 0.071 & 1.27 & 1.65 & --- & --- & --- \\
 & LQRy & 0.040 & 0.71 & 0.93 & --- & --- & --- \\
\addlinespace
(A) Severe mismatch & PID & 3.14 & 818.7 & 554.0 & 44.3 & 644 & 335 \\
 & LQRy & \multicolumn{6}{c}{diverges\textsuperscript{a}} \\
 & LQRy (fixed-gain)\textsuperscript{b} & 0.353 & 2.61 & 1.77 & 0.93 & 0.89 & 0.99 \\
\addlinespace
(B) Longitudinal only & PID & 0.081 & 1.28 & 113.8 & 1.14 & 1.01 & 68.9 \\
 & LQRy & 0.055 & 0.78 & 0.96 & 1.38 & 1.09 & 1.03 \\
\addlinespace
(C) All derivatives $\times1.10$ & PID & 0.073 & 1.28 & 1.65 & 1.03 & 1.01 & 1.00 \\
 & LQRy & 0.032 & 0.72 & 0.93 & 0.82 & 1.01 & 1.00 \\
\addlinespace
(D) Wind and gusts & PID & 0.104 & 2.24 & 2.55 & 1.46 & 1.76 & 1.54 \\
 & LQRy & 0.042 & 0.75 & 5.12 & 1.05 & 1.06 & 5.48 \\
\bottomrule
\multicolumn{8}{l}{\footnotesize\textsuperscript{a}The run terminates
at ${\approx}$55--60\,s, at the first heading-doublet reversal.}\\
\multicolumn{8}{l}{\footnotesize\textsuperscript{b}The earlier
fixed-gain bench build (Section~\ref{sec:hil}); degradation relative
to its own nominal run.}
\end{tabular}
\end{table*}

\subsection{Severe Model Mismatch}
\label{sec:rob-sato}
The perturbed plant replaces the nominal, flight-test-identified
aerodynamic set~\cite{angelo2026} by the complete alternative set
originally distributed with the nonlinear model
implementation~\cite{ferreira2019modeling,repo}, obtained through an
independent identification of the platform. The two sets differ
substantially --- the pitch-damping derivative $C_{m_q}$ by a factor of
4, the elevator effectiveness $C_{m_{\delta_e}}$ by 3.3, and the
lateral--directional set even changes the \emph{sign} of $C_{y_\beta}$
and $C_{y_p}$ --- so this scenario probes well beyond parametric
uncertainty: it is a structurally different aerodynamic description.

Under this mismatch \emph{both} laws lose the mission, each in its own
way. The PID enters a violent but bounded lateral oscillation: the
heading doublet triggers it, the heading error grows to
${\sim}2100^\circ$ (the aircraft flies repeated full circles), and the
altitude sinks by hundreds of meters (Table~\ref{tab:robustness}: RMS
heading error of $819^\circ$, altitude error of 554\,m). The
gain-scheduled LQRy fails harder: the run diverges outright before
60\,s, at the first heading-doublet reversal --- its high-bandwidth
lateral gains amplify the inverted-sign lateral--directional dynamics
instead of containing them: the heading error grows geometrically
across the reversal ($0.6^\circ$ at 51\,s, $17.7^\circ$ by 55\,s)
until the sideslip leaves the physical domain and the simulation
terminates. The contrast within the LQ family itself is the most
instructive result of the scenario: the earlier, softer
\emph{fixed-gain} LQRy iteration --- the build deployed on the bench
of Section~\ref{sec:hil} --- completes this same mission with
\emph{near-nominal} error (degradation of $0.89$--$0.99$,
Table~\ref{tab:robustness} and Fig.~\ref{fig:rob-sato}), so the same
structure spans the full range from unconditional survival to outright
divergence depending only on its tuning. Robustness to severe mismatch
is thus a property of the \emph{tuning}, not of the LQ structure per
se --- a point taken up in Section~\ref{sec:rob-discussion}.
\begin{figure}[!t]\centering
\includegraphics[width=\linewidth]{fig_rob_sato.png}
\caption{Scenario A (severe mismatch): combined mission on the
perturbed plant, nominal-design controllers. The fixed-gain LQRy
(bench build) completes the mission with near-nominal error; the
gain-scheduled LQRy diverges at the first heading reversal
($\times$ marker; inset: 45--58\,s).}\label{fig:rob-sato}\end{figure}

\subsection{Isolating the Origin: Longitudinal Mismatch Only}
\label{sec:rob-satoL}
Scenario B replaces only the \emph{longitudinal} derivatives by the
mismatched set, keeping the nominal lateral--directional model. Both
laws now complete the mission, which identifies the
lateral--directional half of the mismatch --- with its sign inversions
--- as what defeats them both in Scenario~A. The LQRy stays near
nominal in every channel (degradation of 1.03--1.38). The PID
maintains heading and airspeed at nominal levels (degradation of 1.01
and 1.14), but the longitudinal mismatch alone is not benign for it:
the altitude channel degrades by a factor of 69 (RMS 113.8\,m),
a steady drift of the altitude hold. The mechanism is specific: the
$\theta^{ref}$ clamp of $\pm10^\circ$ (Section~\ref{sec:pid}), sized for
the control authority of the nominal model, leaves the altitude loop
without enough authority on the perturbed plant, and the loop drifts
without ever diverging.
\begin{figure}[!t]\centering
\includegraphics[width=\linewidth]{fig_rob_satoL.png}
\caption{Scenario B (longitudinal mismatch only).}
\label{fig:rob-satoL}\end{figure}

\subsection{Plausible Identification Uncertainty}
\label{sec:rob-pert}
Scenario C multiplies all 23 stability and control derivatives by 1.10
--- a uniform $+10\%$ multiplicative perturbation with preserved signs,
representative of a realistic identification error, in the spirit of
the perturbed-derivative robustness campaigns of the fixed-wing UAV
literature~\cite{safwat2019design}. Both controllers are
insensitive to it: every degradation factor lies between 0.82 and 1.03,
and the responses are visually indistinguishable from the nominal runs.
Within plausible identification uncertainty, \emph{both} designs are
robust; the contrast between them only appears under the severe,
sign-changing mismatch of Scenarios A--B.

\subsection{Wind and Gusts}
\label{sec:rob-vento}
Small, low-speed fixed-wing UAVs are disproportionately sensitive to
atmospheric disturbance --- gust velocities are a large fraction of
their airspeed~\cite{mohamed2014fixedwing} --- so no comparison of
flight controllers for this class is complete without a wind scenario.
Scenario D perturbs the nominal plant through the wind inputs
of~\eqref{eq:inputvec}: a 4\,m/s tailwind ramped in over 20--25\,s, two
lateral $1-\cos$ gusts of $\pm3$\,m/s (55--65\,s and 105--115\,s, timed
over the heading doublet), and two vertical $1-\cos$ gusts of
$\pm2$\,m/s (85--95\,s and 165--175\,s, timed over the altitude
doublet). The $1-\cos$ profile is the classical discrete-gust
idealization~\cite{hoblit1988gust}, retained to this day in the gust
conditions of certification regulations~\cite{far25341}; deterministic
discrete gusts (rather than stochastic Dryden or von K\'arm\'an
turbulence~\cite{milf8785c1980,milhdbk1797}) keep the scenario exactly
reproducible run to run, consistent with the deterministic-bench
philosophy of Section~\ref{sec:hil}. Both controllers remain stable
and complete the mission with moderate peak errors. The LQRy shows the
largest single degradation of the scenario --- altitude RMS
$\times5.48$ (5.12\,m) --- while the PID degrades most in heading
under the lateral gusts (RMS $\times1.76$, peak error $12.9^\circ$).
In \emph{absolute} terms, however, the LQRy still tracks heading more
tightly under the gusts (RMS $0.75^\circ$ versus $2.24^\circ$) but
pays in altitude (RMS 5.12\,m versus 2.55\,m). Neither controller
approaches loss of control.
\begin{figure}[!t]\centering
\includegraphics[width=\linewidth]{fig_rob_vento.png}
\caption{Scenario D (wind and gusts): wind profile and closed-loop
responses.}\label{fig:rob-vento}\end{figure}

\subsection{Discussion}
\label{sec:rob-discussion}
The study still exposes \emph{complementary} weaknesses, but
redistributed with respect to the earlier fixed-gain iteration of this
comparison. The severity of the model mismatch now defeats \emph{both}
laws: the cascade PID because a sign change in the lateral derivatives
invalidates its fixed loop shapes, and the gain-scheduled LQRy because
its high-bandwidth lateral gains turn the same sign change into
outright divergence. That an earlier, softer fixed-gain LQRy iteration
tolerated the same scenario is consistent with the gain and phase
margins classically guaranteed for LQ state
feedback~\cite{safonov1977gain} being lost under output feedback in
general~\cite{doyle1978guaranteed}: whatever mismatch tolerance an LQ
output-feedback design exhibits is an empirical property of the
specific tuning rather than a theorem, and aggressive tuning spends
it. The discriminators between the two laws lie elsewhere: under the
longitudinal-only mismatch (B) the PID silently starves an outer loop
through a clamp sized for the nominal plant while the LQRy stays near
nominal, and under wind (D) the LQRy concedes altitude where the PID
holds firm; within plausible uncertainty (C) both are robust.
Robust-control syntheses that shape this trade-off explicitly have
been pursued for small fixed-wing UAVs in our own research
context~\cite{abdulhamid2016robusto}; here the point is narrower:
between the two \emph{standard} laws compared, neither dominates under
uncertainty, and which weakness matters more depends on whether the
dominant threat is model error or atmospheric disturbance.
